% Einheit 4 von 5 – Dreiecke konstruieren (bank/winkel-dreiecke/e4.jsonl, zusammenbau v0.1)
%% TODO Zweigzeile Teil 1: was hier gelernt wird (Satz in Schülersprache aus dem Einheitstitel) – Einheit 4
\einheitenkopf[e4]{Einheit 4 von 5 · Dreiecke konstruieren}
\zweigzeile{ab Kl. 6, je nach Buch bis Kl. 7 (am Gymnasium ab Kl. 5) · P10}
\setcounter{aufgabe}{19}

%% TODO Titel: Ich-kann-Satz fehlt in der Bank (Platzhalter: Kettenname) – Nr. 20
%% TODO Anweisung: ein Satz über der ersten Teilaufgabe fehlt in der Bank – Nr. 20
\begin{aufgabe}{Konstruieren}
%% TODO \rechenplatz{n}: Schrittzahl des Grundfalls fehlt in der Bank (nur bei fester Schreibform, 2.3 b)
\begin{teile}
\teil Gegeben sind a = 6 cm, b = 4 cm, $\gamma = 52^\circ$. Markiere die Stücke in einer Planfigur. Welcher Kongruenzsatz passt? \\ \kreuz{SSS} \\ \kreuz{SWS} \\ \kreuz{WSW} \\ \kreuz{SsW}
\teil Konstruiere das Dreieck ABC mit a = 6 cm, b = 5 cm und c = 7 cm. Beginne mit der Seite c auf dem Strahl.

\winkelstrahl[9][A]
\teil Konstruiere das Dreieck ABC mit b = 5 cm, c = 6 cm und $\alpha = 50^\circ$. Beginne mit der Seite c auf dem Strahl.

\winkelstrahl[9][A]
\teil Konstruiere das Dreieck ABC mit c = 7 cm, $\alpha = 38^\circ$ und $\beta = 71^\circ$. Beginne mit der Seite c auf dem Strahl.

\winkelstrahl[9][A]
\teil Konstruiere das Dreieck ABC mit b = 6 cm, c = 5 cm und $\beta = 64^\circ$. Beginne mit der Seite c auf dem Strahl.

\winkelstrahl[9][A]
\teil Konstruiere ein gleichschenkliges Dreieck ABC mit der Basis c = 6 cm und den Schenkeln a = b = 5 cm.

\winkelstrahl[9][A]
\teil Konstruiere ein rechtwinkliges Dreieck ABC mit dem rechten Winkel bei C und den Katheten a = 4 cm und b = 6 cm.

\winkelstrahl[9][C]
\teil Schreibe eine Konstruktionsbeschreibung in drei bis vier Schritten für das Dreieck ABC mit c = 6 cm, $\alpha = 55^\circ$ und b = 4 cm.
\teil Lässt sich aus den Seiten 4 cm, 5 cm und 10 cm ein Dreieck konstruieren? \\ \kreuz{ja} \\ \kreuz{nein}
\teil Dreieck ABC: AB = 5 cm, BC = 7 cm, CA = 6 cm. Dreieck PQR: PQ = 7 cm, QR = 6 cm, RP = 5 cm. Die Dreiecke sind kongruent. Welche Ecke von PQR gehört zu A, zu B und zu C?
\teil Ein dreieckiges Beet ABC hat die Seite a = 12 m. Was gilt für die Summe der beiden anderen Seiten b und c? \hfill (P10 2020 OS) \\ \kreuz{$b + c > 12$ m} \\ \kreuz{$b + c = 12$ m} \\ \kreuz{$b + c < 12$ m}
\end{teile}
\end{aufgabe}

%% TODO Titel: Ich-kann-Satz fehlt in der Bank (Platzhalter: Kettenname) – Nr. 21
%% TODO Anweisung: ein Satz über der ersten Teilaufgabe fehlt in der Bank – Nr. 21
\begin{aufgabe}{Dreieck skizzieren und beschriften}
\begin{teile}
\teil Beschrifte die Planfigur: Ecken A, B, C gegen den Uhrzeigersinn, dazu die Seiten a, b, c und die Winkel α, β, γ. Markiere farbig die gegebenen Stücke b = 5 cm, c = 6 cm und $\alpha = 38^\circ$.

\dreieck{(0,0)}{(5,0)}{(1.5,3)}{}{}{}{}{}{}
\end{teile}
\end{aufgabe}

%% TODO Titel: Ich-kann-Satz fehlt in der Bank (Platzhalter: Kettenname) – Nr. 22
%% TODO Anweisung: ein Satz über der ersten Teilaufgabe fehlt in der Bank – Nr. 22
\begin{aufgabe}{Konstruieren – Fehler finden, Begründen, Anwendung}
\begin{teile}
\teil Sina sagt: „Mit $\alpha = 52^\circ$, $\beta = 64^\circ$ und $\gamma = 64^\circ$ ist das Dreieck eindeutig festgelegt – das ist der Kongruenzsatz WWW.“ Finde den Fehler.
\teil Warum reicht es nicht, von einem Dreieck nur die drei Winkel zu kennen? Begründe.
\teil Ein dreieckiges Beet soll mit Brettern von 3,5 m, 2,0 m und 1,2 m Länge eingefasst werden. Geht das? Begründe.
\end{teile}
\end{aufgabe}
